Prostorové rozložení míry \ac{AROPE} v~EU27 odhaluje zřetelnou geografickou dichotomii: severozápadní Evropa (CZ, AT, NL, DK, FI) dosahuje hodnot pod \SI{18}{\percent}, zatímco jihovýchodní periferie (RO, BG, GR) překračuje \SI{30}{\percent}.
Pozoruhodná je pozice ČR --- jako jedné z~mála středoevropských zemí se stabilně drží v~dolním kvartilu --- avšak jak ukazuje časová osa (obr.~\ref{fig:arope_timeline_CE}), tato pozice je křehká a~závislá na zachování vysoké zaměstnanosti.
Geografická blízkost zemí s~nízkým \ac{AROPE} a~silným kolektivním vyjednáváním (AT, DE) není náhodná: přeshraniční pracovní mobilita přináší tlak na mzdovou konvergenci zespodu, ale bez domácích institucí \ac{KV} zůstává tento tlak nevyužit.
Česká republika tak paradoxně těží z~blízkosti k~vysokomzdovým ekonomikám jako spotřebitel jejich poptávky po pracovní síle, ale nebuduje institucionální kapacitu, která by umožnila obdobnou mzdovou úroveň doma.

Prostorové rozložení míry \ac{AROPE} v~EU27 odhaluje zřetelnou geografickou dichotomii: severozápadní Evropa (CZ, AT, NL, DK, FI) dosahuje hodnot pod \SI{18}{\percent}, zatímco jihovýchodní periferie (RO, BG, GR) překračuje \SI{30}{\percent}.
Pozoruhodná je pozice ČR --- jako jedné z~mála středoevropských zemí se stabilně drží v~dolním kvartilu --- avšak jak ukazuje časová osa (obr.~\ref{fig:arope_timeline_CE}), tato pozice je křehká a~závislá na zachování vysoké zaměstnanosti.
Geografická blízkost zemí s~nízkým \ac{AROPE} a~silným kolektivním vyjednáváním (AT, DE) není náhodná: přeshraniční pracovní mobilita přináší tlak na mzdovou konvergenci zespodu, ale bez domácích institucí \ac{KV} zůstává tento tlak nevyužit.
Česká republika tak paradoxně těží z~blízkosti k~vysokomzdovým ekonomikám jako spotřebitel jejich poptávky po pracovní síle, ale nebuduje institucionální kapacitu, která by umožnila obdobnou mzdovou úroveň doma.

Prostorové rozložení míry \ac{AROPE} v~EU27 odhaluje zřetelnou geografickou dichotomii: severozápadní Evropa (CZ, AT, NL, DK, FI) dosahuje hodnot pod \SI{18}{\percent}, zatímco jihovýchodní periferie (RO, BG, GR) překračuje \SI{30}{\percent}.
Pozoruhodná je pozice ČR --- jako jedné z~mála středoevropských zemí se stabilně drží v~dolním kvartilu --- avšak jak ukazuje časová osa (obr.~\ref{fig:arope_timeline_CE}), tato pozice je křehká a~závislá na zachování vysoké zaměstnanosti.
Geografická blízkost zemí s~nízkým \ac{AROPE} a~silným kolektivním vyjednáváním (AT, DE) není náhodná: přeshraniční pracovní mobilita přináší tlak na mzdovou konvergenci zespodu, ale bez domácích institucí \ac{KV} zůstává tento tlak nevyužit.
Česká republika tak paradoxně těží z~blízkosti k~vysokomzdovým ekonomikám jako spotřebitel jejich poptávky po pracovní síle, ale nebuduje institucionální kapacitu, která by umožnila obdobnou mzdovou úroveň doma.

Prostorové rozložení míry \ac{AROPE} v~EU27 odhaluje zřetelnou geografickou dichotomii: severozápadní Evropa (CZ, AT, NL, DK, FI) dosahuje hodnot pod \SI{18}{\percent}, zatímco jihovýchodní periferie (RO, BG, GR) překračuje \SI{30}{\percent}.
Pozoruhodná je pozice ČR --- jako jedné z~mála středoevropských zemí se stabilně drží v~dolním kvartilu --- avšak jak ukazuje časová osa (obr.~\ref{fig:arope_timeline_CE}), tato pozice je křehká a~závislá na zachování vysoké zaměstnanosti.
Geografická blízkost zemí s~nízkým \ac{AROPE} a~silným kolektivním vyjednáváním (AT, DE) není náhodná: přeshraniční pracovní mobilita přináší tlak na mzdovou konvergenci zespodu, ale bez domácích institucí \ac{KV} zůstává tento tlak nevyužit.
Česká republika tak paradoxně těží z~blízkosti k~vysokomzdovým ekonomikám jako spotřebitel jejich poptávky po pracovní síle, ale nebuduje institucionální kapacitu, která by umožnila obdobnou mzdovou úroveň doma.

\input{texparts/python/arope_map_2025}



